\section{Einleitung}

\subsection{Problemstellung und Relevanz}

Laut Mitteilung des Statistischen Bundesamtes waren 2024 etwa 15,2 \% aller Kinder und Jugendlichen in Deutschland armutsgefährted \parencite{destatis2025}. Während die Armutsgefährdungsquote relativ zum mittleren Nettoäquivalenzeinkommen berechnet wird und damit eher eine Verteilung misst, wurden auch konkrete Auswirkungen auf die Lebenssituation der Kinder gemessen. So mussten rund 12 \% der Kinder und Jugendlichen auf eine einwöchige Urlaubsreise verzichten und 5 \% konnten aus finanziellen Gründen nicht an regelmäßigen Freizeitaktivitäten teilnehmen \parencite{destatis2025}.

In Deutschland wurden und werden unterschiedliche Maßnahmen ergriffen, um die Armut von Kindern und Jugendlichen zu reduzieren. Neben staatlichen Leistungen wie dem Kindergeld, dem Kinderzuschlag oder dem Bildungs- und Teilhabepaket gibt es Städte und Organisationen, die ihrerseits Initiativen ins Leben rufen. Um diese Maßnahmen kontinuierlich zu verbessern und an die Bedürfnisse der Betroffenen anzupassen, müssen die Auswirkungen der Maßnahmen regelmäßig evaluiert werden. Dabei ist es von besonderer Bedeutung sowohl qualitative als auch quantitative Methoden zu kombinieren, d.h. eine methodische Triangulation durchzuführen. Wird etwa auf quantitativer Ebene die Anzahl in Anspruch genommener Leistungen erhoben, so kann daraufhin eine qualitative Untersuchung nähere Einblicke geben, wie die Leistungen von Betroffenen wahrgenommen wurden oder aus welchen Gründen sie nicht in Anspruch genommen wurden.

\subsection{Ziel der Arbeit und Forschungsfrage}

Ziel dieser Arbeit ist es, die methodische Triangulation als Möglichkeit der Evaluation von Maßnahmen zur Armutsbekämpfung bei Kindern und Jugendlichen zu untersuchen. Dies soll anhand einer konkreten Evaluation geschehen, der Evaluation des Tübinger Präventionskonzeptes gegen Kinderarmut, die im Rahmen einer Kooperation zwischen der Stadt Tübingen und der Evangelischen Hochschule Ludwigsburg durchgeführt wurde \parencite{tuebingen2024}. Mit einem Mixed-Methods-Design wurden in der Evaluation sowohl quantitative als auch qualitative Daten erhoben. Über die Analyse des Erkenntnisgewinnes durch die methodische Triangulation hinaus, soll auch untersucht werden, inwiefern die methodische Triangulation zu höherer Zuverlässigkeit und Validität der Ergebnisse beitragen kann.

Die Forschungsfrage lautet: Welchen Erkenntnisgewinn bietet die Triangulation in der Evaluation von Präventionsmaßnahmen gegen Kinderarmut und inwiefern kann sie zu höherer Zuverlässigkeit und Validität der Ergebnisse beitragen?

\subsection{Aufbau der Arbeit}

Die Arbeit gliedert sich in sieben Kapitel. Nach dieser Einleitung folgt im zweiten Kapitel eine theoretische Einführung in die quantitative und qualitative Forschungsmethodik, das Mixed-Methods-Design, sowie die methodische Triangulation. Im dritten Kapitel wird die ausgewählte Evaluation des Tübinger Präventionskonzeptes gegen Kinderarmut vorgestellt. Im darauffolgenden vierten Kapitel erfolgt eine Analyse der quantitativen und qualitativen Forschungsmethoden. Die Untersuchung der Kombination der quantitativen und qualitativen Forschungsmethoden in einer methodischen Triangulation erfolgt im fünften Kapitel mit einer anschließenden kritischen Bewertung der methodischen Triangualtion im sechsten Kapitel. Das siebte Kapitel fasst die Ergebnisse zusammen, beantwortet die Forschungsfrage und gibt einen Ausblick auf zukünftige Forschungen in diesem Bereich.

% TODO die Motivation stärker Richtung Triangulation lenken, häufig spreche ich von methodischer Triangulation, dabei wäre der umfassenderere Bergiff der Triangulation korrekter